\begin{frame}{지수 감쇠: DAgger 없이 BC만 쓰는 정량적 비용}
  \begin{columns}[T]
    \begin{column}{0.5\textwidth}
      \centering
      \includegraphics[width=\linewidth]{ch12_3_compounding_error.png}
      {\footnotesize $p\in\{0.98,0.99,0.995\}$의 $p^T$ 곡선.}
    \end{column}
    \begin{column}{0.48\textwidth}
      \begin{itemize}
        \item $p=0.995$이면 같은 300스텝에서 \textbf{22\%}까지
          살아남음.
        \item \textbf{1스텝 성공률 0.5\%p}가 긴 episode에서
          ``거의 불가능''과 ``가끔 된다''를 가른다.
        \item DAgger가 ``시연에 없던 상태로 비틀어진'' 경로의
          상태까지 학습 데이터에 포함시켜 \textbf{$p$ 자체를
          높이는} 이유가 바로 이것.
        \item 실습: $p$와 $T$를 바꿔 $p^T$ 곡선을 직접 그려보기.
      \end{itemize}
    \end{column}
  \end{columns}
  \vskip0.5em
  {\footnotesize 그리고 그 위에서 \textbf{RL로 미세조정}하면 그 상한을
  \textbf{넘을 수 있을까} — 다음 실험에서 직접 확인.}
\end{frame}
